
\subsection{Poliedrinės grupės}

Poliedrinės grupės (ang. \emph{polyhedral groups}) yra $\Orth3$ pogrupiai,
kurie veikimu į $R^3$ fiksuoja poliedro (briaunainio) taškus.
Egzistuoja septynios poliedrinės grupės.

\begin{definition}
	\textquote{Chiralinė tetraedro simetrija} yra grupė 
	\begin{equation}
		T = \present{s,t}{s^2=t^3=(st)^3=e}
		\,,
	\end{equation}
	čia generatorius $s$ yra tapatinamas su sukimu kampu $\pi$ aplink ašį einančia per tetraedro briaunos vidurio tašką,
	generatorius $t$ yra tapatinamas su sukimu kampu $2\pi/3$ aplink ašį einančia per tetraedro sienos centrą,
	elementas $st$ yra tapatinamas su sukimu kampu $2\pi/3$ aplink ašį einančia per tetraedro viršūnę.
\end{definition}

\begin{definition}
	\textquote{Chiralinė oktaedro simetrija} yra grupė 
	\begin{equation}
		O = \present{s,t}{s^2=t^3=(st)^4=e}
		\,,
	\end{equation}
	čia generatorius $s$ yra tapatinamas su sukimu kampu $\pi$ aplink ašį einančia per oktaedro briaunos vidurio tašką,
	generatorius $t$ yra tapatinamas su sukimu kampu $2\pi/3$ aplink ašį einančia per oktaedro sienos centrą,
	elementas $st$ yra tapatinamas su sukimu kampu $2\pi/4$ aplink ašį einančia per oktaedro viršūnę.
\end{definition}

\begin{definition}
	\textquote{Chiralinė ikosaedro simetrija} yra grupė 
	\begin{equation}
		I = \present{s,t}{s^2=t^3=(st)^5=e}
		\,,
	\end{equation}
	čia generatorius $s$ yra tapatinamas su sukimu kampu $\pi$ aplink ašį einančia per ikosaedro briaunos vidurio tašką,
	generatorius $t$ yra tapatinamas su sukimu kampu $2\pi/3$ aplink ašį einančia per ikosaedro sienos centrą,
	elementas $st$ yra tapatinamas su sukimu kampu $2\pi/5$ aplink ašį einančia per ikosaedro viršūnę.
\end{definition}

\begin{definition}
	\textquote{Pilnoji tetraedro simetrija} yra grupė 
	\begin{equation}
		T_d = \present{s,t,v}{s^2=t^3=(st)^3=e,\, vsv^{-1}=s^1, vtv^{-1}=t^{-1}}
		\,
	\end{equation}
	čia elementai $s$, $t$ ir $st$ tapatinami taip pat kaip grupėje $T$,
	generatorius $v$ yra tapatinamas su atspindžiu per plokštumą,
	kuriai priklauso tetraedro briauna ir centras.
\end{definition}

\begin{definition}
	\textquote{Pilnoji oktaedro simetrija} yra grupė 
	\begin{equation}
		O_h = \present{s,t}{
			s^2=t^3=(st)^4=i^2=e,\,
			is=si ,\,
			it=ti }
		\,,
	\end{equation}
	čia elementai $s$, $t$ ir $st$ tapatinami taip pat kaip grupėje $O$,
	generatorius $i$ yra tapatinamas su inversija.
\end{definition}

\begin{definition}
	\textquote{Pilnoji ikosaedro simetrija} yra grupė 
	\begin{equation}
		I_h = \present{s,t}{
			s^2=t^3=(st)^5=i^2=e,\,
			is=si ,\,
			it=ti }
		\,,
	\end{equation}
	čia elementai $s$, $t$ ir $st$ tapatinami taip pat kaip grupėje $I$,
	generatorius $i$ yra tapatinamas su inversija.
\end{definition}

\begin{definition}
	\textquote{Piritoedro simetrija} yra grupė
	\begin{equation}
		T_h = \present{s,t}{
			s^2=t^3=(st)^3=i^2=e,\,
			is=si ,\,
			it=ti }
	\end{equation}
	čia elementai $s$, $t$ ir $st$ tapatinami taip pat kaip grupėje $T$,
	generatorius $i$ yra tapatinamas su inversija.
\end{definition} 

Grupės $T$, $O$ ir $I$ yra \emph{chiralinės grupės},
nes jos yra grupės $\SO3$ pogrupiai.
Grupės $T_d$, $O_h$, $I_h$ ir $T_h$ yra \emph{achiralinės grupės},
nes joms priklauso netikrieji pasukimai.

